\subsection{\texorpdfstring{Exactness criteria for the functors
$H_\bullet(P_\bullet\otimes_A M)$.}{Exactness criteria for the functors H.(P. tensor M)}}
\label{subsection:III.7.4}
\label{III.7.4}

\begin{env}[7.4.1]
\label{III.7.4.1}
Let $A$ be a ring, not necessarily commutative, and let
$P_\bullet$ be a complex of \emph{flat} right $A$-modules.  Since the
functor
\[
  M\leadsto P_k\otimes_A M
\]
is then exact in $\cat{Ab}_A$ for every $k$, the
$\partial$-functor
\[
  T_\bullet(M)=H_\bullet(P_\bullet\otimes_A M)
  \tag{7.4.1.1}\label{III.7.4.1.1.eq}
\]
\oldpage[III]{186}
is a \emph{homological functor} from $\cat{Ab}_A$ to
$\cat{Ab}$; it is plainly $A$-linear when $A$ is commutative
\sref{III.7.1.2}, and it commutes with inductive limits.

If $A$ is commutative, then, for every $A$-algebra $B$, the
homological functor $T_\bullet^{(B)}$ \sref{III.7.3.8} is, by
definition, given by
\[
  T_\bullet^{(B)}(N)
  =
  H_\bullet(P_\bullet\otimes_A N_{[\rho]}),
  \tag{7.4.1.2}\label{III.7.4.1.2.eq}
\]
where $\rho:A\to B$ is the homomorphism defining the
$A$-algebra structure on $B$.  Since one may also write $P_\bullet\otimes_A N_{[\rho]}=P_\bullet\otimes_A(B\otimes_B N)_{[\rho]}=(P_\bullet\otimes_A B)\otimes_B N$,
it follows that
\[
  T_\bullet^{(B)}(N)
  =
  H_\bullet(P'_\bullet\otimes_B N)
  \tag{7.4.1.3}\label{III.7.4.1.3.eq}
\]
for every $B$-module $N$, where $P'_\bullet$ is the complex
$P_\bullet\otimes_A B$ of \emph{flat} $B$-modules
\sref[0]{0.6.2.1}.
\end{env}

\begin{proposition}[7.4.2]
\label{III.7.4.2}
Under the general hypotheses of \sref{III.7.4.1}, and for a
given integer $p\in\bb{Z}$, the following conditions are
equivalent:
\begin{enumerate}
  \item[(a)] $T_p$ is left exact \textup{(or, what amounts to the same,}
    $T_{p+1}$ is right exact);
  \item[(b)] $Z'_p(P_\bullet)
=
\Coker(P_{p+1}\longrightarrow P_p)$
    is a flat right $A$-module;
  \item[(c)] there exists a complex $P'_\bullet$ of flat right
    $A$-modules for which the differential
    \[
      d_{p+1}:P'_{p+1}\longrightarrow P'_p
    \]
    is zero, together with an isomorphism of homological
    functors
    \[
      H_\bullet(P_\bullet\otimes_A M)
      \simeq
      H_\bullet(P'_\bullet\otimes_A M).
    \]
\end{enumerate}
\end{proposition}

\begin{proof}
By definition, there is an exact sequence, functorial in $M$,
\[
  0
  \longrightarrow
  T_p(M)
  \longrightarrow
  Z'_p(P_\bullet\otimes M)
  \longrightarrow
  P_{p-1}\otimes M,
\]
where $Z'_p(P_\bullet\otimes M)
=
\Coker(P_{p+1}\otimes M\longrightarrow
P_p\otimes M)
=
Z'_p(P_\bullet)\otimes M$
by the right exactness of the tensor product.  For every
homomorphism $f:M\to N$ one therefore has a commutative diagram
\[
\xymatrix{
  0 \ar[r]
    & T_p(M) \ar[r] \ar[d]_{u}
    & Z'_p(P_\bullet)\otimes M \ar[r] \ar[d]_{v}
    & P_{p-1}\otimes M \ar[d]_{w} \\
  0 \ar[r]
    & T_p(N) \ar[r]
    & Z'_p(P_\bullet)\otimes N \ar[r]
    & P_{p-1}\otimes N
}
\tag{7.4.2.1}\label{III.7.4.2.1.eq}
\]
whose rows are exact.  If $f$ is a monomorphism, so is $w$,
since $P_{p-1}$ is flat; if $T_p$ is left exact, $u$ is also a
monomorphism; one concludes that $v$ is a monomorphism, which
implies that $Z'_p(P_\bullet)$ is flat.  Conversely, if this is
so, $v$ is a monomorphism for every monomorphism $f:M\to N$, so
diagram \eqref{III.7.4.2.1.eq} shows that $u$ is a
monomorphism, and consequently $T_p$ (which is already
semi-exact) is left exact.  Thus (a) and (b) are equivalent.
It is immediate that (c) implies (a), because if
$d_{p+1}:P_{p+1}\to P_p$ is zero and
$0\to M'\to M\to M''\to0$ is an exact sequence of $A$-modules,
the boundary operator in the exact sequence
\[
  H_{p+1}(P_\bullet\otimes M'')
  \xrightarrow{\partial}
  H_p(P_\bullet\otimes M')
  \longrightarrow
  H_p(P_\bullet\otimes M)
\]
\oldpage[III]{187}
is zero by definition (M, IV, 1), so $T_p$ is left exact.  Let
us prove conversely that (b) implies (c).  If
$Z_{p+1}(P_\bullet)=\Ker(P_{p+1}\longrightarrow P_p)$, there
is an exact sequence
\[
  0\longrightarrow Z_{p+1}(P_\bullet)\longrightarrow P_{p+1}
  \longrightarrow Z'_p(P_\bullet)\longrightarrow0
\]
in which $P_{p+1}$ and $Z'_p(P_\bullet)$ are flat, so
$Z_{p+1}(P_\bullet)$ is flat \sref[0]{0.6.1.2}.  We shall take
\[
  P'_i=P_i\quad\text{for }i\ne p\text{ and }i\ne p+1,
  \qquad
  P'_p=Z'_p(P_\bullet)
  \qquad\text{and}\qquad
  P'_{p+1}=Z_{p+1}(P_\bullet);
\]
as differential $d'_i:P'_i\to P'_{i-1}$ we shall take that of
the complex $P_\bullet$ for $i\ne p$ and $i\ne p+1$, $0$ for
$i=p+1$, and for $i=p$ the homomorphism
$Z'_p(P_\bullet)\to P_{p-1}$ deduced from $d_p$ by passage to
the quotient.  Since the $P_i$ are flat, one has
\[
  Z'_i(P_\bullet\otimes M)=Z'_i(P_\bullet)\otimes M,
  \qquad
  Z_i(P_\bullet\otimes M)=Z_i(P_\bullet)\otimes M
  \quad\text{and}\quad
  B_i(P_\bullet\otimes M)=B_i(P_\bullet)\otimes M
\]
(putting $B_i(P_\bullet)=\Im(P_{i+1}\to P_i)$); one immediately
concludes that there are, for every $M$, functorial isomorphisms
$H_i(P_\bullet\otimes M)\xrightarrow{\sim}H_i(P'_\bullet\otimes M)$
for every $i$, and the verification that they form an
isomorphism of $\partial$-functors follows without difficulty
from the definition of $\partial$ (M, IV, 1).
\end{proof}

Note that the conditions of \sref{III.7.4.2} also imply that
$B_p(P_\bullet)$ is \emph{flat}, since there is an exact sequence
$0\longrightarrow B_p(P_\bullet)\longrightarrow P_p
\longrightarrow Z'_p(P_\bullet)\longrightarrow0$
in which $P_p$ and $Z'_p(P_\bullet)$ are flat
\sref[0]{0.6.1.2}.

\begin{corollary}[7.4.3]
\label{III.7.4.3}
Suppose that $A$ is a regular Noetherian ring of dimension
$1$, in other words, a \emph{product of Dedekind rings}
\sref[0]{0.17.1.3} and \sref[0]{0.17.3.7}; for example, $A$
may be a \emph{principal} ideal ring.  Then $T_p$ is left exact if
and only if $T_p(A)$ is a flat $A$-module.  The functor $T_p$
is exact if and only if both $T_p(A)$ and $T_{p-1}(A)$ are
flat $A$-modules.
\end{corollary}

\begin{proof}
Recall that, for a module $M$ over a Dedekind ring, being \emph{flat}
is equivalent to being \emph{torsion-free} \sref[0]{0.6.3.3} and
\sref[0]{0.6.3.4}.  Under the hypotheses of
\sref{III.7.4.3}, every submodule of a flat $A$-module is
therefore flat.

The second assertion follows from the first, since saying that
$T_p$ is exact amounts to saying that $T_p$ and $T_{p-1}$ are
left exact.  To prove the first assertion, observe that there
is an exact sequence
\[
  0
  \longrightarrow H_p(P_\bullet)
  \longrightarrow Z'_p(P_\bullet)
  \longrightarrow B_{p-1}(P_\bullet)
  \longrightarrow 0,
\]
where $B_{p-1}(P_\bullet)$ is \emph{flat}, being a submodule of the
flat $A$-module $P_{p-1}$.  Hence $H_p(P_\bullet)$ is flat if
and only if $Z'_p(P_\bullet)$ is flat
\sref[0]{0.6.1.2}.
\end{proof}

The most important applications of \sref{III.7.4.2} are the
following.

\begin{proposition}[7.4.4]
\label{III.7.4.4}
Let $A$ be a Noetherian ring and $P_\bullet$ a complex of flat
$A$-modules.  Suppose either that the $P_i$ are of finite
type, or that the $H_i(P_\bullet)$ are $A$-modules of finite
type and there exists an $i_0$ such that
$H_i(P_\bullet)=0$ for $i<i_0$.  Let $T_\bullet$ be the
homological functor defined by \eqref{III.7.4.1.1.eq}.  Then
the set $U$ of points $y\in\Spec(A)$ such that $(T_p)_y$
\sref{III.7.1.4} is right exact, respectively left exact or
exact, is open in $\Spec(A)$.
\end{proposition}

\begin{proof}
\oldpage[III]{188}
Under the second hypothesis on $P_\bullet$, one may replace
$P_\bullet$ by a complex $P'_\bullet$ of finite free
$A$-modules such that the functor $H_\bullet(P'_\bullet\otimes_A M)$
is isomorphic to $T_\bullet(M)$ as a $\partial$-functor
\sref[0]{0.11.9.3}.  Thus one may always reduce to the first
hypothesis.  In that case the modules $Z'_i(P_\bullet)$ are of
finite type.  Moreover, it is enough to prove the assertions
concerning left exactness (cf.\ \sref{III.7.4.2}, (a)).

Let $x\in U$.  Since $M\leadsto M_x$ is exact, $\bigl(Z'_p(P_\bullet)\bigr)_x
=
Z'_p\bigl((P_\bullet)_x\bigr)$.
Taking \eqref{III.7.4.1.3.eq} into account, the hypothesis and
\sref{III.7.4.2}, (b), show that
$\bigl(Z'_p(P_\bullet)\bigr)_x$ is a flat $A_x$-module.  It is
therefore \emph{free}, since it is of finite type and $A_x$ is a
Noetherian local ring (Bourbaki,
\emph{Alg\`ebre commutative}, chap.~II, \S~3, no.~2, cor.~2
of prop.~5).  It follows that there is an $f\in A$ for which
$\bigl(Z'_p(P_\bullet)\bigr)_f$ is free over $A_f$
(Bourbaki, \emph{Alg\`ebre commutative}, chap.~II, \S~5,
no.~1, cor.\ of prop.~2).  A fortiori,
$\bigl(Z'_p(P_\bullet)\bigr)_y$ is free over $A_y$ for every
$y\in D(f)$, and \sref{III.7.4.2}, (b), completes the proof.
\end{proof}

\begin{corollary}[7.4.5]
\label{III.7.4.5}
Under the hypotheses of \sref{III.7.4.4}, suppose in addition
that $A$ is an integral domain.  Then the set $U$ of
$x\in\Spec(A)$ such that $(T_p)_x$ is exact is open and
nonempty.
\end{corollary}

\begin{proof}
It is enough to show that $(T_p)_x$ is exact at the \emph{generic}
point $x$ of $\Spec(A)$.  This is immediate, since $A_x$ is a
field and every additive functor on $\cat{Ab}_{A_x}$ is exact.
\end{proof}

\begin{proposition}[7.4.6]
\label{III.7.4.6}
Under the general hypotheses of \sref{III.7.4.4}, conditions
(a), (b), and (c) \textup{of \sref{III.7.4.2}} are also equivalent to:
\begin{enumerate}
  \item[(d)] there exists an $A$-module $Q$ and a functorial
    isomorphism
    \[
      T_p(M)\xrightarrow{\sim}\Hom_A(Q,M).
      \tag{7.4.6.1}\label{III.7.4.6.1.eq}
    \]
\end{enumerate}
Moreover, this property determines the $A$-module $Q$ up to a
unique isomorphism, and $Q$ is of finite type.
\end{proposition}

\begin{proof}
The uniqueness of $Q$ is a special case of the uniqueness of a
representing object for a representable functor
\sref[0]{0.8.1.5}.  It is clear that the right-hand side of
\eqref{III.7.4.6.1.eq} is left exact.

Conversely, to prove the existence of $Q$ when $T_p$ is left
exact, one may first, as in \sref{III.7.4.4}, reduce to the
case where the $P_i$ are flat and of finite type, and hence
(since $A$ is Noetherian) \emph{projective} of finite type (Bourbaki,
\emph{Alg\`ebre commutative}, chap.~II, \S~5, no.~2, cor.\ of
th.~1).  The \emph{dual} $\check{P}_i$ of $P_i$ is then also a
projective $A$-module of finite type, $P_i$ is canonically
isomorphic to the dual of $\check{P}_i$, and the canonical
homomorphism $P_i\otimes_A M
\longrightarrow
\Hom_A(\check{P}_i,M)$
is bijective (Bourbaki, \emph{Alg\`ebre}, chap.~II,
3rd ed., \S~4, no.~2, prop.~2).

By \sref{III.7.4.2}, (c), one may furthermore suppose that
$d_{p+1}:P_{p+1}\to P_p$ is zero.  There is then an exact
sequence
\[
  0
  \longrightarrow T_p(M)
  \xrightarrow{u}
  P_p\otimes_A M
  \xrightarrow{v}
  P_{p-1}\otimes_A M,
\]
where $v=d_p\otimes 1$.  Put $Q'=\Ker(d_p)$, so that
\[
  0
  \longrightarrow Q'
  \xrightarrow{w} P_p
  \xrightarrow{d_p}
  P_{p-1}
\]
is exact.  Transposition gives the exact sequence
\[
  \check{P}_{p-1}
  \xrightarrow{{}^t d_p}
  \check{P}_p
  \xrightarrow{{}^t w}
  \check{Q}'
  \longrightarrow 0.
\]
We claim that $Q=\check{Q}'=\Coker({}^t d_p)$
has the required property.  Indeed, there is an exact sequence
\[
  0
  \longrightarrow \Hom(Q,M)
  \longrightarrow \Hom(\check{P}_p,M)
  \xrightarrow{v'}
  \Hom(\check{P}_{p-1},M),
\]
\oldpage[III]{189}
where $v'=\Hom_A({}^t d_p,1)$.  Under the canonical
identifications
\[
  P_i\otimes_A M\simeq\Hom_A(\check{P}_i,M),
\]
the homomorphism $v'$ identifies with $v=d_p\otimes 1$.
Thus one obtains the required functorial isomorphism
\[
  T_p(M)\isoto\Hom_A(Q,M).
\]
Finally, $Q$ is a quotient of $\check{P}_p$, and is therefore
of finite type.
\end{proof}

\begin{proposition}[7.4.7]
\label{III.7.4.7}
Suppose that the general hypotheses of \sref{III.7.4.4} hold.
For every $A$-module $M$ of finite type:
\begin{enumerate}
  \item[\textup{(i)}] the $T_i(M)$ are $A$-modules of finite type;
  \item[\textup{(ii)}] for every ideal $\mathfrak m$ of $A$, the
    canonical homomorphism
    \[
      (T_i(M))^{\widehat{\,}}
      \longrightarrow
      \varprojlim_n
      T_i\bigl(M\otimes_A(A/\mathfrak m^{n+1})\bigr)
      \tag{7.4.7.1}\label{III.7.4.7.1.eq}
    \]
    is bijective, where the left-hand side denotes the
    separated completion of $T_i(M)$ for the
    $\mathfrak m$-adic topology.
\end{enumerate}
\end{proposition}

\begin{proof}
As in \sref{III.7.4.4}, one reduces first to the case where the
$P_i$ are of finite type.  Since $A$ is Noetherian, the
submodules of the $P_i\otimes_A M$ are of finite type, which
immediately proves (i).  Assertion (ii) follows from the more
general lemma below.
\end{proof}

\begin{lemma}[7.4.7.2]
\label{III.7.4.7.2}
Let $A$ be a Noetherian ring and $u:E\to F$ a homomorphism of
finite-type $A$-modules.  For every finite-type $A$-module
$M$, put
\[
  K(M)=\Ker(u\otimes 1_M),
  \qquad
  C(M)=\Coker(u\otimes 1_M).
\]
Then, for every ideal $\mathfrak m$ of $A$, the canonical
homomorphisms
\[
  (K(M))^{\widehat{\,}}
  \longrightarrow
  \varprojlim_n K(M_n),
  \qquad
  (C(M))^{\widehat{\,}}
  \longrightarrow
  \varprojlim_n C(M_n)
  \tag{7.4.7.3}\label{III.7.4.7.3.eq}
\]
are bijective, where $M_n
=
M\otimes_A(A/\mathfrak m^{n+1})
=
M/\mathfrak m^{n+1}M$.
\end{lemma}

\begin{proof}
The modules $E\otimes_A M$ and $F\otimes_A M$ are of finite
type, and the functor $M\leadsto\widehat{M}$ is \emph{exact} on the
category of finite-type $A$-modules \sref[0]{0.7.3.3}.
Consequently, $(K(M))^{\widehat{\,}}$ and $(C(M))^{\widehat{\,}}$ are,
respectively, the kernel and the cokernel of
\[
  (u\otimes 1)^{\widehat{\,}}:
  (E\otimes_A M)^{\widehat{\,}}
  \longrightarrow
  (F\otimes_A M)^{\widehat{\,}}.
\]
Left exactness of $\varprojlim$ therefore gives $(K(M))^{\widehat{\,}}
=
\varprojlim K(M_n)$.
On the other hand, right exactness of the tensor product gives $C(M_n)
=
C(M)\otimes_A(A/\mathfrak m^{n+1})$,
and hence $(C(M))^{\widehat{\,}}
=
\varprojlim C(M_n)$
by definition.
\end{proof}

\begin{remark}[7.4.8]
\label{III.7.4.8}
Taking \sref{III.6.10.5} and \sref{III.6.10.6} into account,
one sees that, with an additional flatness hypothesis,
\sref{III.7.4.7} recovers the fact that
(4.1.7.1) is an isomorphism, which is the
essential content of the ``first comparison theorem'' for
proper morphisms.  Moreover, the statement applies not only
to a coherent $\sh{O}_X$-module but to a \emph{complex} of
such modules.  It would be interesting to obtain a statement
that contains both \sref{III.7.4.7} and
(4.1.7.1) as special cases.

Notice also that, when the $P_i$ are of finite type, the proof
of \sref{III.7.4.7} does not use their flatness.  It would be
worth examining whether the conclusion of \sref{III.7.4.7}
continues to hold when the $P_i$ are assumed neither flat nor
of finite type, but the $H_i(P_\bullet)$ are assumed to be of
finite type for every $i$ and zero for $i<i_0$.  Can one then
replace $P_\bullet$ by a complex $P'_\bullet$ of finite-type
$A$-modules in such a way that the functors
\[
  H_\bullet(P_\bullet\otimes_A M)
  \quad\hbox{and}\quad
  H_\bullet(P'_\bullet\otimes_A M),
\]
which are no longer homological, are still isomorphic?
\end{remark}
